%----------------------------------------------------------------------------------------
% MACHTEN EN WORTELS. Alle opgaven zijn eigen werk.
%
% Een hoofdstuk kernvaardigheden (ontwerpregels, deel Kernvaardigheden): de theorie beknopt herhaald
% (kaders, klassieke fouten, een paar korte voorbeelden), daarna vooral drill: grote Oefenbundels en Test jezelf.
% Labels van dit hoofdstuk beginnen met h1-. Alleen in pakket 5: geen \alleenvoor.
%----------------------------------------------------------------------------------------

\mychapter{1}{Machten en wortels}

\lkr{Herhaling van Reële getallen (3de jaar, LPD 5-7). Nodig in Tweedegraadsfuncties (de discriminant, oplossingen als $1 \pm \sqrt3$), Goniometrie (exacte waarden), Vectoren en Analytische meetkunde (lengtes en afstanden). Bundels 1-5 oefenen telkens één soort, 6-8 mengen, 9 gaat over klassieke fouten en 10 is uitdagend. Laat leerlingen kiezen: wie een bundel vlot maakt, springt naar de volgende.}

\section{Wat je moet kennen}

\begin{theorie}
\begin{minipage}[t]{.47\linewidth}
\vspace{0pt}
\textbf{\textcolor{accent}{Machten}} \quad ($a, b \neq 0$, $m, n$ geheel)

\vspace{1mm}
$\formule{a^m \cdot a^n = a^{m+n}}$ \qquad $\formule{\dfrac{a^m}{a^n} = a^{m-n}}$

\vspace{1mm}
$\formule{(a^m)^n = a^{m \cdot n}}$ \qquad $\formule{(ab)^n = a^nb^n}$

\vspace{1mm}
$\formule{a^0 = 1}$ \qquad $\formule{a^{-n} = \dfrac{1}{a^n}}$
\end{minipage}\hfill
\begin{minipage}[t]{.47\linewidth}
\vspace{0pt}
\textbf{\textcolor{accent}{Vierkantswortels}} \quad ($a, b \geqslant 0$)

\vspace{1mm}
$\formule{\sqrt{ab} = \sqrt a\,\sqrt b}$ \qquad $\formule{\sqrt{\dfrac ab} = \dfrac{\sqrt a}{\sqrt b}}$

\vspace{1mm}
$\formule{\bigl(\sqrt a\bigr)^2 = a}$ \qquad $\formule{\sqrt{a^2} = \abs{a}}$ \ (elke $a$)

\vspace{1mm}
\textbf{\textcolor{accent}{Wetenschappelijke notatie:}} $a \cdot 10^n$ met $1 \leqslant a < 10$
\end{minipage}
\end{theorie}

\textbf{Klassieke fouten.} Er is geen regel voor een som: $(a + b)^2 \neq a^2 + b^2$ en $\sqrt{9 + 16} = 5 \neq \sqrt9 + \sqrt{16}$. De macht gaat voor op het minteken: $-3^2 = -9$, maar $(-3)^2 = 9$. Een negatieve exponent maakt een getal niet negatief: $2^{-3} = \tfrac18$. Wortels tel je alleen op als het dezelfde wortel is: $2\sqrt3 + 5\sqrt3 = 7\sqrt3$.

\voorbeeld{}{Vereenvoudig $\dfrac{(2a^3b^{-1})^2}{4a^{-1}b^3}$ en schrijf zonder negatieve exponenten.\label{vb:h1-macht}}{%
\begin{stappen}
\frac{4a^6b^{-2}}{4a^{-1}b^3} = a^{6-(-1)}\,b^{-2-3} = \frac{a^7}{b^5} & Eerst de macht van het product, dan per grondtal.
\end{stappen}}

\voorbeeld{}{Vereenvoudig: \quad a. $\sqrt{75} - 2\sqrt{12}$ \qquad b. $\dfrac{8}{\sqrt5 - 1}$\label{vb:h1-wortel}}{%
\begin{stappen}
5\sqrt3 - 2 \cdot 2\sqrt3 = \sqrt3 & a. Grootste kwadraatfactor: $75 = 25 \cdot 3$, $12 = 4 \cdot 3$. \\
\frac{8(\sqrt5 + 1)}{(\sqrt5 - 1)(\sqrt5 + 1)} = \frac{8(\sqrt5 + 1)}{4} = 2\sqrt5 + 2 & b. Maal de toegevoegde $\sqrt5 + 1$: de noemer wordt $5 - 1$.
\end{stappen}}

%========================================================================================
\opdrachtenkop

\begin{oefeningen}
\needspace{10\baselineskip}
\oef{h1-o1}

\needspace{10\baselineskip}
\oef{h1-o2}

\needspace{10\baselineskip}
\oef{h1-o3}

\needspace{10\baselineskip}
\oef{h1-o4}

\needspace{10\baselineskip}
\oef{h1-o5}

\needspace{10\baselineskip}
\oef{h1-o6}

\needspace{10\baselineskip}
\oef{h1-o7}

\needspace{10\baselineskip}
\oef{h1-o8}

\needspace{10\baselineskip}
\oef{h1-o9}

\needspace{10\baselineskip}
\oef{h1-o10}
\end{oefeningen}

%----------------------------------------------------------------------------------------
\needspace{14\baselineskip}
\opdrachtsectie{Test jezelf}

\begin{oefeningen}
\oef{h1-t1}

\oef{h1-t2}

\oef{h1-t3}

\oef{h1-t4}

\oef{h1-t5}
\end{oefeningen}

\herhaalschema{%
  \fv{oef:h1-t1}{vb:h1-macht}{oef:h1-o2}%
  \fzv{oef:h1-t2}{oef:h1-o3}%
  \fv{oef:h1-t3}{vb:h1-wortel}{oef:h1-o4}%
  \fv{oef:h1-t4}{vb:h1-wortel}{oef:h1-o5}%
  \fzv{oef:h1-t5}{oef:h1-o9}}

\ifzelfstudie
\section*{Maak je eigen samenvatting}
Sluit dit boek en schrijf in je schrift: de rekenregels voor machten en vierkantswortels, de wetenschappelijke notatie, hoe je een noemer wortelvrij maakt, en drie klassieke fouten. Vergelijk daarna met het kader.
\fi
